\subsection{The scalar components at all twenty-three lambencies}
\label{sec:all-lambency-classification}
The character criterion gives a finite, exact classification of every
untwined component.  We take the independent indices $1\leq r<m$; the
remaining indices are their negatives or are identically zero.

\begin{proposition}[All scalar umbral shadows]\label{prop:all-scalar-shadows}
Among the $303$ components $S^X_r$, with $X$ running through the twenty-three
Niemeier root systems and $1\leq r<m(X)$, exactly $92$ vanish identically,
$39$ have one primitive Dirichlet character in
\eqref{eq:primitive-shadow-expansion}, and $172$ have at least two.
Every one of those $172$ components has the unconditional off-line zeros
of Theorem~\ref{thm:shadow-zero-classification}.
The single-character and zero components are specified in the table below;
all components not listed in either category are mixed.

Of the $39$ single-character components, $21$ have a monomial
$P_\psi(z)$, ten have a monomial times $1+p^{-z}$, and eight have a monomial
times $1+p^{1-z}$, with $p=2$ or $3$.
For these last eight, the additional zeros are
\begin{equation}\label{eq:scalar-shadow-boundary-zeros}
 s=1+\frac{(2j+1)\pi i}{2\log p},\qquad j\in\mathbb Z.
\end{equation}
They are present unconditionally and cannot be cancelled by a Dirichlet
$L$-function, which is nonvanishing at $\Re(2s-1)=1$.
The eight cases are
\[
 \begin{array}{c|c|c}
 X&r&p\\\hline
 A_{11}D_7E_6&4,8&2\\
 E_6^4&4,8&2\\
 A_{17}E_7&9&3\\
 D_{10}E_7^2&3,9,15&3.
 \end{array}
\]
Consequently, under the Dirichlet Riemann hypothesis, every zero of any of
the $39$ single-character components with $\Re s>1/2$ lies on
$\Re s=3/4$ or on the explicitly described line $\Re s=1$.
\end{proposition}

\begin{center}
\small
\begin{tabular}{l r l r r}
$X$&$m$&single-character indices&zero count&mixed count\\\hline
$A_1^{24}$&2&1&0&0\\
$A_2^{12}$&3&1,2&0&0\\
$A_3^8$&4&2&0&2\\
$A_4^6$&5&none&0&4\\
$A_5^4D_4$&6&2,3,4&0&2\\
$A_6^4$&7&none&0&6\\
$A_7^2D_5^2$&8&2,4,6&0&4\\
$A_8^3$&9&3,6&0&6\\
$A_9^2D_6$&10&5&0&8\\
$A_{11}D_7E_6$&12&2,4,6,8,10&0&6\\
$A_{12}^2$&13&none&0&12\\
$A_{15}D_9$&16&4,8,12&0&12\\
$A_{17}E_7$&18&6,9,12&0&14\\
$A_{24}$&25&none&0&24\\
$D_4^6$&6&1,3,5&2&0\\
$D_6^4$&10&5&4&4\\
$D_8^3$&14&7&6&6\\
$D_{10}E_7^2$&18&3,9,15&8&6\\
$D_{12}^2$&22&11&10&10\\
$D_{16}E_8$&30&5,15,25&14&12\\
$D_{24}$&46&23&22&22\\
$E_6^4$&12&4,8&5&4\\
$E_8^3$&30&none&21&8\\\hline
Total&&39&92&172
\end{tabular}
\end{center}
For the seven root systems beginning with $D$, precisely the even indices
vanish.  For $E_6^4$ the zero indices are $2,3,6,9,10$.
For $E_8^3$ the nonzero indices are $1,7,11,13,17,19,23,29$.
Thus the table specifies every index, not only the counts.

\begin{proof}
We describe a finite certificate using exact arithmetic.  For $d\mid m$,
the defining matrix is
\[
 \Omega_m(d)_{r,t}=\begin{cases}
 1,&r+t\equiv0\pmod{2d},\quad r-t\equiv0\pmod{2m/d},\\
 0,&\text{otherwise}.
 \end{cases}
\]
The ADE matrices are, respectively,
\[
 \begin{array}{c|c}
 A_{m-1}&\Omega_m(1)\\
 D_{m/2+1}&\Omega_m(1)+\Omega_m(m/2)\\
 E_6&\Omega_{12}(1)+\Omega_{12}(4)+\Omega_{12}(6)\\
 E_7&\Omega_{18}(1)+\Omega_{18}(6)+\Omega_{18}(9)\\
 E_8&\Omega_{30}(1)+\Omega_{30}(6)+\Omega_{30}(10)+\Omega_{30}(15).
 \end{array}
\]
These are Cheng--Duncan--Harvey's defining matrices
\cite[Table~5 and Sect.~3.3]{CDH2014}; a union of root systems means
adding their matrices.  For each row put
$c_r(n)=\Omega^X_{r,n}-\Omega^X_{r,-n}$, and apply
\eqref{eq:character-stratum-coefficients} followed by
\eqref{eq:imprimitive-shadow}.

The computation enumerates the characters using the cyclic unit groups of
odd prime powers and
$(\mathbb Z/2^a\mathbb Z)^\times=\langle-1\rangle\times\langle5\rangle$
for $a\geq3$.  All sums are evaluated in the exact ring
$\mathbb Q[x]/(\Phi_e(x))$, where $e$ is the exponent of the relevant unit
group.  The conductor is the smallest divisor $f$ of the modulus for which
the character is trivial on the kernel of reduction modulo $f$.
After collecting equal primitive characters, every resulting polynomial
is expanded back into its coefficient sequence.  Equality with $c_r$ is
checked at every residue of its full period $2m$; the reconstructed terms
have periods dividing $2m$, so this check proves each identity.
The complete tables and the program \emph{Shadow character classification}
are in the workbench.

The table then follows by counting the nonzero primitive constituents.
For the single-character cases the exact polynomials are scalar multiples
of $d^{-z}$, $d^{-z}(1+p^{-z})$, or $d^{-z}(1+p^{1-z})$.
To show the last possibility explicitly, the two components with
$X=A_{11}D_7E_6$, $r=4,8$, have
\[
 D_{c_r}(z)=2\,4^{-z}(1+2^{1-z})L(z,\chi_{-3}),
\]
and the two with $X=E_6^4$, $r=4,8$, have twice this expression.
The component $X=A_{17}E_7$, $r=9$, has
\[
 D_{c_9}(z)=3^{-z}(1+3^{1-z})L(z,\chi_{-4});
\]
the components $r=3,15$ of $D_{10}E_7^2$ have this same expression,
and its component $r=9$ has twice it.  Solving $p^{1-z}=-1$ with
$z=2s-1$ gives \eqref{eq:scalar-shadow-boundary-zeros}.
The zero theorem for the other $172$ components is
Theorem~\ref{thm:shadow-zero-classification}.
\end{proof}
